% Options for packages loaded elsewhere
\PassOptionsToPackage{unicode}{hyperref}
\PassOptionsToPackage{hyphens}{url}
\documentclass[
  10pt,
  a4paper,
]{ltjsarticle}
\usepackage{xcolor}
\usepackage[margin=22mm]{geometry}
\usepackage{amsmath,amssymb}
\setcounter{secnumdepth}{-\maxdimen} % remove section numbering
\usepackage{iftex}
\ifPDFTeX
  \usepackage[T1]{fontenc}
  \usepackage[utf8]{inputenc}
  \usepackage{textcomp} % provide euro and other symbols
\else % if luatex or xetex
  \usepackage{unicode-math} % this also loads fontspec
  \defaultfontfeatures{Scale=MatchLowercase}
  \defaultfontfeatures[\rmfamily]{Ligatures=TeX,Scale=1}
\fi
\usepackage{lmodern}
\ifPDFTeX\else
  % xetex/luatex font selection
  \setmainfont[]{Hiragino Mincho ProN}
  \setsansfont[]{Hiragino Sans}
  \setmonofont[]{Menlo}
\fi
% Use upquote if available, for straight quotes in verbatim environments
\IfFileExists{upquote.sty}{\usepackage{upquote}}{}
\IfFileExists{microtype.sty}{% use microtype if available
  \usepackage[]{microtype}
  \UseMicrotypeSet[protrusion]{basicmath} % disable protrusion for tt fonts
}{}
\makeatletter
\@ifundefined{KOMAClassName}{% if non-KOMA class
  \IfFileExists{parskip.sty}{%
    \usepackage{parskip}
  }{% else
    \setlength{\parindent}{0pt}
    \setlength{\parskip}{6pt plus 2pt minus 1pt}}
}{% if KOMA class
  \KOMAoptions{parskip=half}}
\makeatother
\usepackage{longtable,booktabs,array}
\newcounter{none} % for unnumbered tables
\usepackage{calc} % for calculating minipage widths
% Correct order of tables after \paragraph or \subparagraph
\usepackage{etoolbox}
\makeatletter
\patchcmd\longtable{\par}{\if@noskipsec\mbox{}\fi\par}{}{}
\makeatother
% Allow footnotes in longtable head/foot
\IfFileExists{footnotehyper.sty}{\usepackage{footnotehyper}}{\usepackage{footnote}}
\makesavenoteenv{longtable}
\usepackage{graphicx}
\makeatletter
\newsavebox\pandoc@box
\newcommand*\pandocbounded[1]{% scales image to fit in text height/width
  \sbox\pandoc@box{#1}%
  \Gscale@div\@tempa{\textheight}{\dimexpr\ht\pandoc@box+\dp\pandoc@box\relax}%
  \Gscale@div\@tempb{\linewidth}{\wd\pandoc@box}%
  \ifdim\@tempb\p@<\@tempa\p@\let\@tempa\@tempb\fi% select the smaller of both
  \ifdim\@tempa\p@<\p@\scalebox{\@tempa}{\usebox\pandoc@box}%
  \else\usebox{\pandoc@box}%
  \fi%
}
% Set default figure placement to htbp
\def\fps@figure{htbp}
\makeatother
\setlength{\emergencystretch}{3em} % prevent overfull lines
\providecommand{\tightlist}{%
  \setlength{\itemsep}{0pt}\setlength{\parskip}{0pt}}
\usepackage{bookmark}
\IfFileExists{xurl.sty}{\usepackage{xurl}}{} % add URL line breaks if available
\urlstyle{same}
\hypersetup{
  pdftitle={自己無撞着な関係波閉鎖系におけるインフレーション的急拡大の機構------第 2 版：計算条件の訂正と N=3〜16 の再検証},
  pdfauthor={木原 範昭（WF System Co., Ltd.）　ORCID 0009-0004-6753-4020},
  hidelinks,
  pdfcreator={LaTeX via pandoc}}

\title{自己無撞着な関係波閉鎖系におけるインフレーション的急拡大の機構------第
2 版：計算条件の訂正と N=3〜16 の再検証}
\author{木原 範昭（WF System Co., Ltd.）　ORCID 0009-0004-6753-4020}
\date{2026年8月30日　Version DOI 10.5281/zenodo.22176949}

\begin{document}
\maketitle

{
\setcounter{tocdepth}{2}
\tableofcontents
}
\textbf{著者:} 木原 範昭（WF System Co., Ltd.）　\textbf{日付:}
2026-08-30 \textbf{Version DOI:}
\href{https://doi.org/10.5281/zenodo.22176949}{10.5281/zenodo.22176949}　\textbf{Concept
DOI:}
\href{https://doi.org/10.5281/zenodo.22112008}{10.5281/zenodo.22112008}
\textbf{訂正対象:} 第 1 版 Version DOI
\href{https://doi.org/10.5281/zenodo.22112009}{10.5281/zenodo.22112009}（2026-08-27
公開。以下「v1」）

\begin{center}\rule{0.5\linewidth}{0.5pt}\end{center}

\subsection{要旨}\label{ux8981ux65e8}

本稿は v1 の訂正版である。v1
の引用する全プログラムを原本まで遡って再実行・監査した結果、v1
の数値結果は公開プログラムから無変更で再現する一方、そのプログラムに理論と食い違う計算条件が
3 点あった：(1) 初期化 \texttt{make\_parent}
と相互作用の内部に残っていた\textbf{隠れた振幅正規化}、(2)
時間発展に用いていた Cayley
変換（\textbf{凍結生成子の有理近似}。固有位相応答が
\(2\arctan(\gamma\sigma)\) に歪む）、(3)
乱数から作った初期状態が修正後の力学では\textbf{自己無撞着でない}こと。本稿はこの
3 点を、(1) 正規化の全面除去と振幅込み相互作用
\(K_{ij}=\mathrm{Im}(\bar z_i z_j)\)、(2)
各刻みで生成子を凍結した厳密直交の指数回転 \(\exp(\Delta K)\)（連続流
\(dv/d\tau=K(v)v\)
の、ノルムと閉塞を厳密に保存する一次積分器。以下「線形回転」）、(3)
厳密な自己無撞着解による初期化（等モジュラー制約あり／なしの 2
系統）、に置き換えて \(N=3\)〜16 を再検証した。

主張は 3 つである。

\begin{enumerate}
\def\labelenumi{\arabic{enumi}.}
\tightlist
\item
  \textbf{ゼロ閉塞 \(\sum_m z_m^2=0\) は公理ではなく定理である。}
  生成子が実反対称で、状態がその非零固有モード（自己無撞着）でありさえすれば、正規化・振幅分布・\(N\)・親の生成法に依らず閉塞は成立する（証明
  5 行。4 生成法 54 親と旧 4 系統 56 親の step 0 で
  \(|\sum z^2|/H\le 10^{-13}\)、40000 step の走行中も
  \(\le 5\times10^{-13}\)）。v1 §6・§24
  の導出はこの意味で正しく、本稿はその射程（何にも依らないこと）を確定する。
\item
  \textbf{複素シンプレックスは幾何に制約を与えない。}
  実シンプレックスでは 2
  乗距離が半正定値条件（Schoenberg）を満たす必要があるが、複素対称行列は常に
  Takagi 分解を持つため、\textbf{任意の}複素 2 乗距離の組が \(N-1\)
  次元複素空間に厳密に埋め込める。乱数状態 1400 個（\(N=3\)〜16 ×
  100）の全てが誤差 \(10^{-15}\) で埋め込まれ rank \(N-1\)
  を持つのに対し、実の読み方では \(N\ge6\) で 100 個中 0
  個しか埋め込めない。複素シンプレックスは状態から各関係の符号
  \(z_e\to-z_e\)
  を忘れた像（\(v/(\mathbb Z_2)^M\)）であり、選択原理ではない。閉塞が可能なのは複素化で正値性を捨てたからである。
\item
  \textbf{計算条件の変更はインフレーション的発展に少なくない影響を与えたが、インフレーション的発展そのものは再現した。}
  修正後の力学でも、自己無撞着で不安定な初期状態からは計算精度の底（倍精度で
  \(10^{-32}\)。多倍長ではさらに深く、底は物理でなく精度と親残差が決める）から指数成長し飽和する。本稿で確認したのは、自己無撞着な相対平衡のうち線形不安定なものが任意に小さい摂動（親残差・丸め誤差）を指数増幅して非線形飽和すること------すなわち線形不安定性とその飽和------であり、摂動源を含む「開始機構」ではない。成長率は
  v1 の \(0.17\)〜\(0.25\)/step から \(0.006\)〜\(0.022\)/step へ 10〜30
  倍遅くなり、飽和後の状態は v1 の等分配（PR/M
  \(=1\)）ではなく少数の関係への局在（PR/M
  \(=0.05\)〜0.31）になる。高対称な初期状態の系列（1-因子分解／距離クラス構成、等モジュラーとその振幅変形）では
  \(N=6,8,\dots,16\) が全て飽和し \(N=5,7,\dots,15\) が全て床（\(N=3,4\)
  はともに床）という明瞭な奇偶非対称が現れたが、対称性のない乱数均衡親では奇数
  \(N\) にも不安定状態があるので、これは \(N\)
  の偶奇だけで決まる普遍則ではない。4 生成法 × \(N=3\)〜16 の 54
  走行について、走行前に固定した共回転 1
  刻み線形化行列（離散写像のヤコビアン）による予測は 53/54
  で実測分類と一致し、飽和した 25 走行の成長率は予測と 0.997〜1.008
  で一致した。
\end{enumerate}

本稿は新しい解釈を加えない。v1 の各主張の維持／修正／撤回を §9
の表で明示し、v1
で暗黙に仮定されていた「自己無撞着性が物理状態を一意に選ぶ」という命題が成り立たないことを記録して閉じる。

\begin{center}\rule{0.5\linewidth}{0.5pt}\end{center}

\subsection{1. v1
からの訂正の位置づけ}\label{v1-ux304bux3089ux306eux8a02ux6b63ux306eux4f4dux7f6eux3065ux3051}

\subsubsection{1.1 訂正の連鎖}\label{ux8a02ux6b63ux306eux9023ux9396}

先行研究 {[}K8{]} は外部 seed
を除去しても急拡大が起きることを示し、「初期化 \texttt{make\_parent}
は零二乗閉塞条件を前提としている」と記した。v1 §2.1
はコード監査によりこれを訂正し、\texttt{make\_parent}
は閉塞を拘束として課しておらず、自己無撞着固定点探索の出力として閉塞が現れると述べた。この訂正は正しい。

本稿はさらに一段遡る。v1 §4 は \texttt{make\_parent}
を「監査済み」としたが、監査は 2
点を見落としていた：\texttt{make\_parent} 内の
\texttt{v\ =\ v/‖v‖}（親の振幅正規化）と、相互作用を位相だけで組む
\(K_{ij}=\sin(\theta_j-\theta_i)\)（相互作用内部の隠れた振幅正規化）である。加えて時間発展の
Cayley 変換は凍結生成子に対する有理近似（固有位相応答
\(2\arctan(\gamma\sigma)\)）であり、凍結生成子の厳密な指数回転（応答
\(\Delta\sigma\)）に置き換えるべきものであった。v1
の数値主張はこれらの計算条件の下での結果である。

\subsubsection{1.2
本稿の書き方}\label{ux672cux7a3fux306eux66f8ux304dux65b9}

本稿の全ての主張は次の 4 種類のいずれかに分類し、各所で明記する。

\begin{itemize}
\tightlist
\item
  \textbf{証明済み}：公理と定義から計算だけで従うもの。
\item
  \textbf{数値事実}：どの点（\(N\)・生成法・走行数）で何が測られたかを添えて述べるもの。
\item
  \textbf{未証明命題}：数学的に定式化できるが証明も反例も無いもの。
\item
  \textbf{証明に必要な補題}：未証明命題を定理にするために示すべき有限個の事実。
\end{itemize}

「〜を強く支持する」「〜と考えられる」という表現は用いない。

また本稿で「インフレーション的」と呼ぶのは、模型内で
\(H_\perp/H_{\rm total}\)
が指数増幅して非線形飽和する現象を指し、宇宙論的インフレーションとの同一性を主張するものではない（v1
§25--26 の宇宙論的類似は本稿では扱わない）。

\begin{center}\rule{0.5\linewidth}{0.5pt}\end{center}

\subsection{2. 系の最小入力（v1 §3・§5
の継承）}\label{ux7cfbux306eux6700ux5c0fux5165ux529bv1-35-ux306eux7d99ux627f}

\subsubsection{2.1
状態と相互作用}\label{ux72b6ux614bux3068ux76f8ux4e92ux4f5cux7528}

\(N\) 個の存在の全二体関係 \(M=\binom N2\) のそれぞれに複素数の波
\(z_e=a_e+ib_e\) を置く。状態ベクトルは
\(v=(z_e)\in\mathbb C^M\)、そのノルム
\(\|v\|^2=\sum_e|z_e|^2=H_{\rm total}\)
である。関係グラフの構造パラメータは \(N\) のみで、\(M\)
は導出量である（v1 §3 の「独立パラメータは \(N\)
のみ」は構造についての記述であり、後述の刻み \(\Delta\)
と全体振幅------\(K(cv)=|c|^2K(v)\)
により回転速度に効く------は別に指定する）。

相互作用は、点を共有する 2 つの関係 \(e,f\) の間だけに働き、その強さは

\[
K_{ef}=A_{ef}\,\mathrm{Im}(\bar z_e z_f)=A_{ef}\,(a_eb_f-b_ea_f)
\]

である（\(A\) は関係の線グラフの隣接行列）。\(K\) は実反対称 \(K^T=-K\)
で、\(K(cz)=|c|^2K(z)\) を満たす（振幅に比例する）。

\subsubsection{2.2 発展}\label{ux767aux5c55}

発展則は連続流

\[
\frac{dv}{d\tau}=K(v)\,v
\]

である。\(K\) は実反対称なので流れは \(\|v\|^2\) と \(v^Tv=\sum z_e^2\)
を保存し、\(H_{\rm int}=\tfrac12\sum_{\{e,f\}}(\mathrm{Im}\,\bar z_ez_f)^2\)（和は隣接する非順序対
\(\{e,f\}\) を各 1 回だけ数える）に対するハミルトン流
\(\dot z_e=2i\,\partial H_{\rm int}/\partial\bar z_e=(Kz)_e\)
でもある（本稿は \(\dot z=+2i\,\partial H/\partial\bar z\)
の符号規約を採る。\(\dot z=-2i\,\partial H/\partial\bar z\) の流儀では
\(-H_{\rm int}\)
がハミルトニアンになる）。数値計算では各刻みで生成子を凍結した指数写像

\[
v_{t+1}=\exp(\Delta\,K(v_t))\,v_t,\qquad \Delta=\frac{2\pi}{L}
\]

で進める。\(\exp(\Delta K)\) は実直交行列なので \(\|v\|^2\) と \(v^Tv\)
は離散写像でも厳密に保存される（v1 §16 の定理は Cayley
でも指数写像でも同じ）が、\(H_{\rm int}\) は厳密には保存されず
\(O(\Delta)\)
で漂う（§10）。本稿で「線形回転」と呼ぶのはこの凍結生成子の指数写像のことで、旧コードの
Cayley 変換（凍結生成子の有理写像）に対する呼称である。どちらも凍結した
\(K\) に対しては線形作用素であり、状態 \(v\)
に対する写像としてはどちらも非線形である。刻みの通し番号は処理ステップであり、流れの径数
\(\tau\) も物理的な時間ではない（v1 §12、および
§8.2：時間に相当する量は状態自身の位相進みとして別に定義される）。

\subsubsection{2.3
自己無撞着状態}\label{ux81eaux5df1ux7121ux649eux7740ux72b6ux614b}

\(v\) が自分自身の生成子の非零固有モード

\[
iK(v)\,v=\mu v,\qquad \mu\neq0
\]

であるとき自己無撞着と呼ぶ。このとき \(v\) は形を変えずに角速度 \(\mu\)
で一様回転する（相対平衡）。\(v=a+ib\) と書けば
\(Ka=\mu b\)、\(Kb=-\mu a\)
であり、実部と虚部は実反対称作用の固有モード内部に現れる回転対である（v1
§5：実数の位相と実反対称作用から複素回転構造が出る）。

最小入力は v1 §3 と同じく

\[
N+\text{完全二体関係}+\text{自己無撞着固定点条件}
\]

である。本稿で新たに明らかになるのは、この最小入力が状態を\textbf{一意に定めない}ことである（§7.4・§11）。

\begin{center}\rule{0.5\linewidth}{0.5pt}\end{center}

\subsection{3. 3
つの変更と、その切り分け}\label{ux3064ux306eux5909ux66f4ux3068ux305dux306eux5207ux308aux5206ux3051}

\subsubsection{3.1 変更の内容}\label{ux5909ux66f4ux306eux5185ux5bb9}

v1 が引用する 15
パッケージの原本プログラム（\texttt{run\_n\_scaling\_lowrank\_v1\_*.py}
系エンジン）と、本稿の修正後プログラムの差は次の 3 点に整理される。

{\def\LTcaptype{none} % do not increment counter
\begin{longtable}[]{@{}
  >{\raggedright\arraybackslash}p{(\linewidth - 4\tabcolsep) * \real{0.3333}}
  >{\raggedright\arraybackslash}p{(\linewidth - 4\tabcolsep) * \real{0.3333}}
  >{\raggedright\arraybackslash}p{(\linewidth - 4\tabcolsep) * \real{0.3333}}@{}}
\toprule\noalign{}
\begin{minipage}[b]{\linewidth}\raggedright
変更
\end{minipage} & \begin{minipage}[b]{\linewidth}\raggedright
旧（v1）
\end{minipage} & \begin{minipage}[b]{\linewidth}\raggedright
新（本稿）
\end{minipage} \\
\midrule\noalign{}
\endhead
\bottomrule\noalign{}
\endlastfoot
\textbf{1. 隠れた振幅正規化の除去} & (a) \texttt{make\_parent} 内で親を
\(v\mapsto v/\|v\|\) に正規化。(b)
相互作用を位相のみで組む：\(K_{ef}=\sin(\theta_f-\theta_e)\)（\(K(cz)=K(z)\)、振幅を捨てる）。(c)
一部パッケージで初期状態に外部 seed を加えて正規化 & (a) 除去。(b)
振幅込み \(K_{ef}=\mathrm{Im}(\bar z_ez_f)\)。(c)
除去（\(Z_0=v\)、種なし） \\
\textbf{2. Cayley 有理写像から凍結生成子の指数写像へ} & Cayley 変換
\((I-\gamma K)^{-1}(I+\gamma K)\)、\(\gamma=\tan(\pi/144)\)（凍結 \(K\)
の有理関数、固有位相応答は \(2\arctan(\gamma\sigma)\)）。一部に
\(K/\sigma_{\max}\) 正規化 & 凍結 \(K\) の厳密直交指数回転
\(\exp(\Delta K)\)、\(\Delta=2\pi/L\)（固有位相応答は
\(\Delta\sigma\)）。\(K/\sigma\) 枝は廃止。どちらも凍結 \(K\)
に対しては線形、状態に対しては非線形（§2.2） \\
\textbf{3. 初期化の変更} & 乱数位相から位相のみ \(K\)
の固有モード反復で得た状態（修正後の力学の固定点ではない：振幅込み \(K\)
に対する残差 \(0.058\)〜\(0.007\)） &
厳密な自己無撞着解。3-1：等モジュラー（全関係の \(|z_e|\)
が等しい）制約あり。3-2：制約なし（非等モジュラー） \\
\end{longtable}
}

\subsubsection{3.2 v1
パッケージの無変更再実行（数値事実）}\label{v1-ux30d1ux30c3ux30b1ux30fcux30b8ux306eux7121ux5909ux66f4ux518dux5b9fux884cux6570ux5024ux4e8bux5b9f}

v1 が引用する 15 パッケージをパス変更 4
箇所以外は無変更で再実行し、保存データ・図と突合した（\texttt{論文v1\_全再現テスト\_20260828}）。主要数値主張------Floquet
乗数 \(\mu_1=1.090086569\)（9 桁）、onset--残差則の傾き 11.616、成長率
0.172513、厳密保存 \(4.4\times10^{-16}\)、等分配、閉包探索の exact
部分------は再現した。再現しないのは丸め誤差の指数増幅で軌道が分かれる後半の量（秩序化の到達
step、\(H_\parallel/H_\perp\)
の最終分配、モジュライの位相値、出発深さの桁数）である。\(N=5\) の
\(3+3+2+2\) 時間分離（2627/4923
step）を生成したプログラムは公開パッケージに含まれていなかった。

\subsubsection{3.3
三者比較・四者比較による切り分け（数値事実）}\label{ux4e09ux8005ux6bd4ux8f03ux56dbux8005ux6bd4ux8f03ux306bux3088ux308bux5207ux308aux5206ux3051ux6570ux5024ux4e8bux5b9f}

同じ 15
パッケージに変更を段階的に適用して再実行した（\texttt{論文v1\_全プログラム修正版\_20260828}、\(L=144\)、5000
step）。

{\def\LTcaptype{none} % do not increment counter
\begin{longtable}[]{@{}
  >{\raggedright\arraybackslash}p{(\linewidth - 4\tabcolsep) * \real{0.3333}}
  >{\raggedright\arraybackslash}p{(\linewidth - 4\tabcolsep) * \real{0.3333}}
  >{\raggedright\arraybackslash}p{(\linewidth - 4\tabcolsep) * \real{0.3333}}@{}}
\toprule\noalign{}
\begin{minipage}[b]{\linewidth}\raggedright
系
\end{minipage} & \begin{minipage}[b]{\linewidth}\raggedright
適用した変更
\end{minipage} & \begin{minipage}[b]{\linewidth}\raggedright
結果
\end{minipage} \\
\midrule\noalign{}
\endhead
\bottomrule\noalign{}
\endlastfoot
原本 & なし & 潜伏 100〜600 step → 指数成長 0.17〜0.25/step →
等分配（\(|z|^2\to1/M\)、PR/M \(=1\)）、\(\sigma=N-1\)、Floquet
1.0901 \\
baseline & 変更 2 と変更 1(a)(c)（相互作用は位相のみのまま） &
\textbf{原本を再現}：成長率 0.172（\(N=5\)）、Floquet
1.0903、等分配、\(N=5\) の 13 個・12 通り閉包も再現 \\
fixed & 変更 1(b) も適用、初期化は旧のまま & 親が固定点でないため step 1
から弾道的に離脱、指数域なし、\(H_\perp/H\to0.93\)〜0.999、局在（\(|z|^2\)
相対幅 \(\approx2\)）、スペクトルエントロピー 0.87 \\
fixed\_equimodular & 変更 1・2・3-1 &
潜伏が戻り、丸め床から指数成長。\(N=5\)：5000 step
内では床（\(4.6\times10^{-19}\)）、\(N=6,7\)：onset 4152/4316、成長率
0.0059/0.0063、\(N=8\)〜15：床、\(N=16\)：\(1.3\times10^{-17}\)。Floquet（\(N=5\)）1.0001 \\
\end{longtable}
}

この表から次が確定する。

\begin{itemize}
\tightlist
\item
  \textbf{変更 2（Cayley→指数写像）と変更
  1(a)(c)（親と初期状態の正規化）は v1 の現象の原因ではない}（baseline
  が原本を再現）。
\item
  \textbf{v1 固有の像------数百 step の潜伏と \(0.17\)〜\(0.25\)/step
  の急拡大、等分配、\(N=5\) の \(3+3+2+2\)------は変更
  1(b)、すなわち位相のみの相互作用に依存していた}。位相のみ \(K\)
  は消えた波が満額の作用を出す相互作用であり、等分配はその帰結であった。一方、\(\sigma=N-1\)
  はスケール不変形 \(\mu=-(N-1)r^2\) として振幅込み \(K\)
  でも成り立ち（§6.1）、インフレーション的成長そのものも振幅込み \(K\)
  で残る（§7）。
\item
  変更 1(b) を入れると旧初期状態は固定点でなくなる（変更 3
  が必要になる理由）。fixed
  の弾道離脱は親の残差が作る項であり、不安定性ではない（§8.1）。
\end{itemize}

fixed の \(N=6,7,10,11\) は変更 1(a)
の副作用で親探索の収束判定（単位ノルム前提の残差式）が壊れ、対照として無効である。fixed\_equimodular
ではこの問題は起きない（残差 \(10^{-11}\)〜\(10^{-13}\)）。

\pandocbounded{\includegraphics[keepaspectratio,alt={図 1}]{figs/figA_four_way_Hperp_N3_N16.png}}

\textbf{図 1.} 四者比較の \(H_\perp/H_{\rm total}\)（\(L=144\)、5000
step、差し引き読出し・床
\(10^{-16}\)）、\(N=3\)〜16。灰：原本（Cayley・位相のみ
\(K\)・親正規化）、青：baseline（線形回転・位相のみ
\(K\)）、赤：fixed（振幅込み
\(K\)・旧親）、緑：fixed\_equimodular（振幅込み
\(K\)・等モジュラー自己無撞着親）。baseline は原本と重なり、fixed は
step 1 から弾道離脱し、fixed\_equimodular は \(N=6,7\)（および \(N=10\)
の後半）だけが 5000 step 内に指数成長を見せる。

\subsubsection{3.4
統一プロトコル}\label{ux7d71ux4e00ux30d7ux30edux30c8ux30b3ux30eb}

§5 以降の新しい走行は全て次のプロトコルで行う：振幅込み
\(K_{ef}=\mathrm{Im}(\bar z_ez_f)\)、\(\exp((2\pi/124)K)\)、正規化なし、種なし
\(Z_0=v\)、40000 step、\(H_\perp\) は親平面（\(a,b\) の張る実 2
次元）に直交する成分の直接計算
\(\|Z-p(p\cdot Z)-q(q\cdot Z)\|^2\)（丸め床 \(10^{-32}\)）。刻み
\(L\)（§3.3 は 144、§7 は
124）は、連続極限の軌道には時間の再径数化としてしか効かない（v1
§12）が、有限刻みの離散計算では \(O(\Delta^2)\)
の積分器依存を残す（§8.3：中立な親の見かけの成長率 \(\Delta^2/8\) は
\(L\) で変わる）。

\begin{center}\rule{0.5\linewidth}{0.5pt}\end{center}

\subsection{4. 主張
1：ゼロ閉塞は定理である}\label{ux4e3bux5f35-1ux30bcux30edux9589ux585eux306fux5b9aux7406ux3067ux3042ux308b}

\subsubsection{4.1
定理と証明（証明済み）}\label{ux5b9aux7406ux3068ux8a3cux660eux8a3cux660eux6e08ux307f}

\textbf{定理.} \(K\) が実反対称行列で、\(v=a+ib\) が
\(iKv=\mu v\)、\(\mu\ne0\) を満たすなら
\(v^Tv=\sum_e z_e^2=0\)、すなわち \(a^Ta=b^Tb\) かつ \(a^Tb=0\)。

\textbf{証明.} \(iK(a+ib)=\mu a+i\mu b\) の実部・虚部から
\(Ka=\mu b\)、\(Kb=-\mu a\)。よって
\(\mu\,a^Ta=-a^TKb\)、\(\mu\,b^Tb=b^TKa=-a^TKb\)（反対称性）、したがって
\(\mu(a^Ta-b^Tb)=0\)。また
\(\mu\,a^Tb=a^TKa=0\)（反対称行列の二次形式は零）。\(\mu\ne0\)
なので主張が従う。□

証明に使ったのは「\(K\) が実」「\(K\) が反対称」「\(v\)
が非零固有モード」の 3
つだけである。正規化・振幅分布・\(N\)・隣接構造・親の生成法は現れない。同じことは固定した
\(K\) の固有ベクトル \(v\) と \(\bar v\)（固有値 \(\mu\) と
\(-\mu\)）の直交性 \(\bar v^\dagger v=v^Tv=0\) としても言える（v1 §6
の形）。

\textbf{系（力学による保存）.} \(\exp(\Delta K)\) は実直交なので
\(v^Tv\) は保存される。閉塞は初期状態で成立すれば全 step で成立する。

\textbf{注意（射程）.} この定理は位相のみの \(K\)（v1
の旧エンジン）でも成立する。したがって閉塞の成立は振幅込み相互作用の正しさの証拠にはならない。閉塞の起源は、実反対称生成子の非零複素固有モード構造（\(K\)
が実、\(K\) が反対称、\(\mu\ne0\) の 3 条件）にある。

\subsubsection{4.2
実測（数値事実）}\label{ux5b9fux6e2cux6570ux5024ux4e8bux5b9f}

\textbf{旧 4
系統（2026-08-29）}：原本／baseline／fixed／fixed\_equimodular ×
\(N=3\)〜16 の 56 親で、step 0 の \(|\sum z^2|/H\) は最大
\(1.0\times10^{-13}\)（fixed\_equimodular、親残差 \(10^{-11}\)
相当）、他は \(\le4.4\times10^{-15}\)。振幅が等分配でも局在（\(|z|^2\)
相対幅
1.3）でも、正規化あり（\(\|v\|=1\)）でも無し（\(0.45\)〜\(1.91\)）でも成立。

\textbf{本稿の 4 生成法（§7.1）}：54 親で
\(\le8\times10^{-14}\)（make\_parent）、\(\le2.5\times10^{-16}\)（他 3
法）。走行中の \(|Z^TZ|/H\) は 40000 step で
\(\le5.1\times10^{-13}\)（丸め蓄積のみ）。

\pandocbounded{\includegraphics[keepaspectratio,alt={図 2}]{figs/fig1_closure_step0.png}}

\textbf{図 2.} step 0 の \(|\sum z^2|/H\)、\(N=3\)〜16。塗り記号は本稿の
4 生成法（54 親）、× は 2026-08-29 の参照 4 系統（56
親）。生成法・正規化の有無・振幅分布に依らず全て \(10^{-13}\)
以下（make\_parent 系は親残差 \(10^{-11}\)
相当、手作り・非等モジュラー・乱数均衡は \(10^{-16}\) 台）。

\subsubsection{4.3 v1 との関係}\label{v1-ux3068ux306eux95a2ux4fc2}

v1 §6 は同じ定理を導出し、§24
で「閉塞は公理から導出可能な定理」と位置づけていた。本稿はこれを取り消さない。訂正するのは、v1
では「自己無撞着な親から閉塞が出る」ことを主に旧エンジンの親で確認していた点で、本稿は定理が生成法・正規化・振幅分布のいずれにも依らないことを
110
親の実装で数値検証した（証明済みの定理をサンプルで確定する必要はなく、検証は実装が定理の前提を満たしていることの確認である）。

\(U^n=I\) については v1 §24 のとおり本定理からは従わない（\(S^1\)
軌道のコンパクト性 \(\not\Rightarrow\)
有限回帰）。本稿でも未導出のままである。

\begin{center}\rule{0.5\linewidth}{0.5pt}\end{center}

\subsection{5. 主張
2：複素シンプレックスは制約を与えない}\label{ux4e3bux5f35-2ux8907ux7d20ux30b7ux30f3ux30d7ux30ecux30c3ux30afux30b9ux306fux5236ux7d04ux3092ux4e0eux3048ux306aux3044}

\subsubsection{5.1
定理（証明済み）}\label{ux5b9aux7406ux8a3cux660eux6e08ux307f}

各関係の 2 乗距離を \(d^2_{ij}=z_{ij}^2\)
と読み、\(D^2=(d^2_{ij})\)、\(J=I-\mathbf 1\mathbf 1^T/N\)、\(B=-\tfrac12JD^2J\)
とする（v1 §17 と同じ）。

\begin{itemize}
\tightlist
\item
  \textbf{実の場合}：\(d^2_{ij}\) が実で、\(N\) 点の実配置
  \(x_i\in\mathbb R^{N-1}\) が \(|x_i-x_j|^2=d^2_{ij}\)
  を満たすための必要十分条件は \(B\succeq0\)（Schoenberg
  {[}E3{]}、Cayley--Menger 条件と同値）。これは不等式の拘束で、\(d^2\)
  の空間の中の錐を切り出す。
\item
  \textbf{複素の場合}：\(B\) は複素対称行列。任意の複素対称行列は
  Autonne--Takagi 分解 \(B=U\Sigma U^T\)（\(U\) ユニタリ、\(\Sigma\ge0\)
  対角）を持つ {[}E2{]}。\(X=U\Sigma^{1/2}\) と置けば
  \(B=XX^T\)。\(D^2\) は対角零の対称行列なので
  \(B_{ii}+B_{jj}-2B_{ij}=d^2_{ij}\) が恒等的に成り立ち、\(X\) の行
  \(x_i\) について
  \((x_i-x_j)\cdot(x_i-x_j)=x_i\cdot x_i+x_j\cdot x_j-2x_i\cdot x_j=B_{ii}+B_{jj}-2B_{ij}=d^2_{ij}\)（双線形形式、複素共役なし）。さらに
  \(B\mathbf 1=0\)（\(J\mathbf1=0\)）なので非零特異値の列だけを残した
  \(X\) は \(X^T\mathbf1=0\)、すなわち重心零で rank \(X\le N-1\) の
  \(\mathbb C^{N-1}\)
  埋め込みである。正定値条件は消え、\textbf{任意の}複素 \(d^2\)
  が埋め込める。□
\end{itemize}

したがって \(\mathbb C^M\) の任意の状態 \(v\) は（rank \(<N-1\)
の退化を許した）複素シンプレックスを与える。ただし辞書 \(d^2=z^2\)
は各関係の符号 \(z_e\to-z_e\) を忘れるので、状態から形への写像は 1:1
ではない：零成分がなければ 1 つの形に \(2^M\)
個の状態（符号枝）が対応し、形が表すのは \(v/(\mathbb Z_2)^M\)
である。符号枝は現在の力学では厳密な離散対称性で結ばれる：\(S=\mathrm{diag}(s_e)\)、\(s_e=\pm1\)
に対して \(K(Sv)=SK(v)S\)（\(K_{ef}\) に \(s_es_f\)
が掛かる）、したがって \(\Phi(Sv)=e^{\Delta SK(v)S}Sv=S\,\Phi(v)\)
であり、符号枝どうしの軌道は共役である。形の側の自由度は \(N(N-1)\)
複素座標 \(-(N-1)\) 平行移動 \(-(N-1)(N-2)/2\) 複素直交群
\(O(N-1,\mathbb C)\) \(=M\) で、符号を除いた状態の自由度と一致する。

\textbf{系.} 実の読み方 \(d^2=|z|^2\ge0\) では閉塞 \(\sum d^2=0\)
は全零以外で不可能である。閉塞が可能なのは、複素化が正値性という拘束を除去したからである。複素シンプレックスは公理でも定理でも選択でもなく、\(d^2:=z^2\)
という辞書（規約）のもとで状態の符号を忘れた像である。

\subsubsection{5.2
実測（数値事実）}\label{ux5b9fux6e2cux6570ux5024ux4e8bux5b9f-1}

自己無撞着でも閉塞でもない乱数状態（各成分の実部・虚部が標準正規、\(N=3\)〜16
× 100
個）について（\texttt{v2補完実験\_…/program/pass2\_embed\_random.py}）：

{\def\LTcaptype{none} % do not increment counter
\begin{longtable}[]{@{}
  >{\raggedright\arraybackslash}p{(\linewidth - 6\tabcolsep) * \real{0.2500}}
  >{\raggedright\arraybackslash}p{(\linewidth - 6\tabcolsep) * \real{0.2500}}
  >{\raggedright\arraybackslash}p{(\linewidth - 6\tabcolsep) * \real{0.2500}}
  >{\raggedright\arraybackslash}p{(\linewidth - 6\tabcolsep) * \real{0.2500}}@{}}
\toprule\noalign{}
\begin{minipage}[b]{\linewidth}\raggedright
\(N\)
\end{minipage} & \begin{minipage}[b]{\linewidth}\raggedright
複素：埋め込み誤差 \(\max|(x_i-x_j)^2-d^2_{ij}|/\max|d^2|\)
\end{minipage} & \begin{minipage}[b]{\linewidth}\raggedright
複素：rank \(=N-1\)
\end{minipage} & \begin{minipage}[b]{\linewidth}\raggedright
実（\(d^2=|z|^2\)）：\(B\succeq0\) で埋め込める個数 /100
\end{minipage} \\
\midrule\noalign{}
\endhead
\bottomrule\noalign{}
\endlastfoot
3 & \(1.6\times10^{-15}\) & 100/100 & 63 \\
4 & \(1.3\times10^{-15}\) & 100/100 & 12 \\
5 & \(1.6\times10^{-15}\) & 100/100 & 2 \\
6〜16 & \(\le2.7\times10^{-15}\) & 100/100 & 0 \\
\end{longtable}
}

\pandocbounded{\includegraphics[keepaspectratio,alt={図 3}]{figs/fig2_embed_random.png}}

\textbf{図 3.} 乱数状態 100 個 × \(N=3\)〜16
の埋め込み。左：複素（\(d^2=z^2\)）の最大相対誤差------全て \(10^{-15}\)
台で厳密に埋め込まれ rank \(=N-1\)。右：実の読み（\(d^2=|z|^2\)）で
\(B\succeq0\) となり埋め込める個数------\(N\ge6\) で 0。

\subsubsection{5.3 v1
の幾何定理の位置づけ}\label{v1-ux306eux5e7eux4f55ux5b9aux7406ux306eux4f4dux7f6eux3065ux3051}

v1 §17 の「全頂点の star 二乗閉塞 \(\Leftrightarrow\) 全頂点が
\(x_i\cdot x_i=0\) 上にある」は \(B\)
の代数だけで成り立つ定理であり、維持する。ただし名称は「光円錐」から\textbf{複素ヌル錐}に改める：複素二次形式の零集合には時間的／空間的の区別がない。\(N=5\)
の等モジュラー状態のように \(d^2\) が全て実で \(\mathbb R^{2,2}\)
に実埋め込みできる場合だけ光円錐と呼べるが、それは等モジュラー点の性質であって
\(N=5\) の性質ではない（§7.1 の非等モジュラー \(N=5\) では \(d^2\)
は複素になる）。

本稿の結論として、複素シンプレックスは\textbf{結果を描く言語}であって\textbf{状態を選ぶ原理}ではない。v1
が §18
で「準安定状態＝等モジュラー・ヌル複素シンプレックス」と述べたとき、幾何が状態を選んでいたのではなく、位相のみの相互作用がそこへ収束させていた（§3.3）。

\begin{center}\rule{0.5\linewidth}{0.5pt}\end{center}

\subsection{6. 主張 3
の準備：自己無撞着な初期状態の作り方}\label{ux4e3bux5f35-3-ux306eux6e96ux5099ux81eaux5df1ux7121ux649eux7740ux306aux521dux671fux72b6ux614bux306eux4f5cux308aux65b9}

変更 3 に用いた初期状態（以下「親」）の 4 生成法を定義する。以下
\(S_i=\sum_{j\ne i}z_{ij}^2\)（頂点 \(i\) の局所和。局所閉塞は
\(S_i=0\)）、\(W_i=\sum_{j\ne i}|z_{ij}|^2\)（頂点重み）と書く。

\subsubsection{6.1
自己無撞着の頂点形（証明済み）}\label{ux81eaux5df1ux7121ux649eux7740ux306eux9802ux70b9ux5f62ux8a3cux660eux6e08ux307f}

\(\mathrm{Im}(\bar z_ez_f)=(\bar z_ez_f-z_e\bar z_f)/2i\) を代入すると

\[
(iKv)_e=\tfrac12\bigl[\bar z_e(Az^2)_e-z_e(A|z|^2)_e\bigr],
\]

\(e=(i,j)\) の隣接和は
\((Az^2)_e=S_i+S_j-2z_e^2\)、\((A|z|^2)_e=W_i+W_j-2|z_e|^2\)
なので、自己無撞着 \(iKv=\mu v\) は各関係について

\[
S_i+S_j=(2\mu+W_i+W_j)\,\frac{z_e^2}{|z_e|^2}
\]

と同値である。\textbf{系}：\(S_i=0\)（全頂点）かつ
\(W_i=W\)（全頂点）ならば自己無撞着で \(\mu=-W\)。逆に自己無撞着かつ
\(S\equiv0\) なら \(W_i\equiv-\mu\)。等モジュラー状態では
\(W_i=(N-1)r^2\)、したがって \(\mu=-(N-1)r^2\)（v1 の \(\sigma=N-1\)
のスケール不変形。\(N=3\) は \(-\tfrac32 r^2\)）。

\subsubsection{6.2 4 生成法}\label{ux751fux6210ux6cd5}

{\def\LTcaptype{none} % do not increment counter
\begin{longtable}[]{@{}
  >{\raggedright\arraybackslash}p{(\linewidth - 6\tabcolsep) * \real{0.2500}}
  >{\raggedright\arraybackslash}p{(\linewidth - 6\tabcolsep) * \real{0.2500}}
  >{\raggedright\arraybackslash}p{(\linewidth - 6\tabcolsep) * \real{0.2500}}
  >{\raggedright\arraybackslash}p{(\linewidth - 6\tabcolsep) * \real{0.2500}}@{}}
\toprule\noalign{}
\begin{minipage}[b]{\linewidth}\raggedright
記号
\end{minipage} & \begin{minipage}[b]{\linewidth}\raggedright
生成法
\end{minipage} & \begin{minipage}[b]{\linewidth}\raggedright
制約
\end{minipage} & \begin{minipage}[b]{\linewidth}\raggedright
構成
\end{minipage} \\
\midrule\noalign{}
\endhead
\bottomrule\noalign{}
\endlastfoot
mp & make\_parent 等モジュラー（3-1） & 等モジュラー & 3
段階：旧アルゴリズム（位相のみ \(K\) の固有モード反復、正規化なし）で
rank-2 の状態 → 位相のみ \(K\) の時間発展で \(|z|^2\) の相対幅
\(<10^{-9}\) まで進める → 振幅込み \(K\) の固有モードへの混合反復で残差
\(<10^{-10}\) まで磨く。乱数は rng(40260721+1000N) \\
hm & 手作り等モジュラー（3-1） & 等モジュラー & 偶数 \(N\)：完全グラフの
1-因子分解（円法）の色 \(c\) に位相 \(\theta=c\pi/(N-1)\)。奇数
\(N\)：距離クラス \(d\) に
\(\theta=(d-1)\pi/q\)、\(q=(N-1)/2\)。\(N=3\)：\(0/60/120°\)。一般和則
\(\sum_f r_f^2\sin^2\phi_{ef}=-\mu\)、\(\sum_f r_f^2\sin2\phi_{ef}=0\)
を厳密に満たす（残差 \(\le1.2\times10^{-16}\)） \\
ne & 非等モジュラー（3-2） & なし & hm と同じ位相配置で、クラス \(c\)
の振幅二乗を \(a_c=\bar r^2(1+0.6\cos(4\pi c/q))\)
とした状態（自己無撞着条件は
\(\sum_ca_c\omega^c=0\)、\(\omega=e^{2\pi i/q}\)。\(q\ge4\)
で成立）。\(q\le3\) の \(N=3,4,5,7\)
では等モジュラー点を通る自己無撞着解の連続族の上の別の点（\(N=5,7\) は
\(S=0\) を保つ分枝、\(N=3,4\) は \(S\ne0\)） \\
rb & 乱数均衡（3-2） & なし & 乱数初期値から \(\{S_i=0,\ W_i=W_0\}\) を
Newton で解いた状態（§6.1
の系により自己無撞着）。位相クラスも対称性もない。各 \(N\) 1
個、rng(100+N) \\
\end{longtable}
}

各 \(N\) で \(\|v\|\) を mp
親に揃える（スケールは回転速度にしか効かない）。受け入れ検査（54/54
合格）：残差
\(\le6.4\times10^{-11}\)（mp）／\(\le1.7\times10^{-16}\)（他）、\(|\sum z^2|/H\le8\times10^{-14}\)、\(\mu\ne0\)。局所閉塞は
\(N=3\)（不可能）と ne の \(N=4\) を除いて成立。rank は全て
\(N-1\)（\(N=3\) の hm は 1、mp・ne は 2）。

非等モジュラーな自己無撞着解が存在すること（ne, rb）自体が、v1 §18
の等モジュラーが自己無撞着の必要条件ではないことを示す。等モジュラー点を通る自己無撞着解集合の次元は
\(N\ge5\) で \(N^2-4N+1\)（\(=2M-(3N-1)\)、\(\{S=0,W=\text{const}\}\)
の拘束数と一致）と数値で測られている（\texttt{複素シンプレックス\_重心閉塞\_非等モジュラー族\_20260830}。等モジュラー点近傍の解集合が
\(\{S=0,W=\text{const}\}\) に一致することは、2 つの有限な rank
条件を厳密演算で確認すれば陰関数定理で定理になる：\textbf{証明に必要な補題}）。

\subsubsection{6.3 走行前予測}\label{ux8d70ux884cux524dux4e88ux6e2c}

各親について 1 刻み写像 \(\Phi(z)=\exp(\Delta K(z))z\)
の実ヤコビアンを中心差分で求め、親の回転
\(e^{-i\phi}\)（\(\phi=\Delta\mu\)）を戻した共回転 1 刻み線形化行列
\(G=R(+\phi)D\Phi(v)\)（回転座標系での固定点の接写像。先行記録では共回転モノドロミーと呼んだ）
の最大固有値絶対値 \(\rho\) から \(\lambda_f=2\ln\rho\)（\(H_\perp\) の
1
刻み成長率）を計算し、走行前に固定した（\texttt{results/parents\_predictions.csv}）。判定規則：\(\rho-1>10^{-3}\)
なら「不安定（40000 step
内に飽和）」、以下なら「中立（床）」。既知の床（hm の \(N=5,7\)）でも
\(\rho-1\approx10^{-4}\) が出る（§8.3 の刻み由来項。\(N=4\) は
\(10^{-7}\)）ため、この閾値は既知の床と飽和の実走行 6
種類（\(N=4\)〜9）で 2026-08-30 に較正したものである。

\begin{center}\rule{0.5\linewidth}{0.5pt}\end{center}

\subsection{\texorpdfstring{7. 主張
3：計算条件変更後のインフレーション的発展、\(N=3\)〜16}{7. 主張 3：計算条件変更後のインフレーション的発展、N=3〜16}}\label{ux4e3bux5f35-3ux8a08ux7b97ux6761ux4ef6ux5909ux66f4ux5f8cux306eux30a4ux30f3ux30d5ux30ecux30fcux30b7ux30e7ux30f3ux7684ux767aux5c55n316}

\subsubsection{\texorpdfstring{7.1
結果行列（数値事実。\texttt{v2補完実験\_…/results/matrix\_N\_by\_method.md}）}{7.1 結果行列（数値事実。v2補完実験\_\ldots/results/matrix\_N\_by\_method.md）}}\label{ux7d50ux679cux884cux5217ux6570ux5024ux4e8bux5b9fv2ux88dcux5b8cux5b9fux9a13_resultsmatrix_n_by_method.md}

床＝40000 step で \(\max H_\perp/H<10^{-10}\)。飽和＝\(H_\perp/H\ge0.5\)
に到達。\(t_{50}\)＝到達 step。\(\lambda\)＝\(\ln H_\perp\)
の指数窓（\(10^{-10}<H_\perp<10^{-3}\)）の傾き /step。括弧内は走行前予測
\(\lambda_f\)。

{\def\LTcaptype{none} % do not increment counter
\begin{longtable}[]{@{}
  >{\raggedright\arraybackslash}p{(\linewidth - 8\tabcolsep) * \real{0.2000}}
  >{\raggedright\arraybackslash}p{(\linewidth - 8\tabcolsep) * \real{0.2000}}
  >{\raggedright\arraybackslash}p{(\linewidth - 8\tabcolsep) * \real{0.2000}}
  >{\raggedright\arraybackslash}p{(\linewidth - 8\tabcolsep) * \real{0.2000}}
  >{\raggedright\arraybackslash}p{(\linewidth - 8\tabcolsep) * \real{0.2000}}@{}}
\toprule\noalign{}
\begin{minipage}[b]{\linewidth}\raggedright
\(N\)
\end{minipage} & \begin{minipage}[b]{\linewidth}\raggedright
mp（等モジュラー）
\end{minipage} & \begin{minipage}[b]{\linewidth}\raggedright
hm（等モジュラー）
\end{minipage} & \begin{minipage}[b]{\linewidth}\raggedright
ne（非等モジュラー）
\end{minipage} & \begin{minipage}[b]{\linewidth}\raggedright
rb（乱数均衡）
\end{minipage} \\
\midrule\noalign{}
\endhead
\bottomrule\noalign{}
\endlastfoot
3 & 床 & 床 & 床 & --- \\
4 & 床 & 床 & 床 & --- \\
5 & \textbf{飽和} \(t_{50}=5029\), \(\lambda=0.0089\) (0.0089) & 床
(0.0002) & \textbf{飽和} 5971, 0.0113 (0.0114) & \textbf{飽和} 5322,
0.0126 (0.0126) \\
6 & \textbf{飽和} 2933, 0.0151 (0.0151) & \textbf{飽和} 4338, 0.0160
(0.0160) & \textbf{飽和} 4395, 0.0148 (0.0148) & \textbf{飽和} 6115,
0.0103 (0.0103) \\
7 & \textbf{飽和} 6564, 0.0070 (0.0070) & 床 (0.0002) & 床 (0.0003) &
\textbf{飽和} 6553, 0.0097 (0.0097) \\
8 & 床・成長中 \(\lambda=0.0005\) (0.0005) & \textbf{飽和} 3864, 0.0177
(0.0177) & \textbf{飽和} 4031, 0.0172 (0.0172) & \textbf{飽和} 9345,
0.0067 (0.0066) \\
9 & 床・成長中 0.0007 (0.0007) & 床 (0.0003) & 床 (0.0003) &
\textbf{飽和} 6505, 0.0098 (0.0098) \\
10 & 床 (0.0003) & \textbf{飽和} 3468, 0.0198 (0.0198) & \textbf{飽和}
3610, 0.0193 (0.0193) & \textbf{飽和} 6522, 0.0100 (0.0100) \\
11 & 床 (0.0003) & 床 (0.0003) & 床 (0.0003) & \textbf{飽和} 6814,
0.0094 (0.0094) \\
12 & 床 (0.0004) & \textbf{飽和} 3391, 0.0202 (0.0202) & \textbf{飽和}
3371, 0.0198 (0.0198) & \textbf{飽和} 11535, 0.0057 (0.0057) \\
13 & 床 (0.0003) & 床 (0.0003) & 床 (0.0003) & \textbf{飽和} 30809,
0.0020 (0.0020) \\
14 & 床 (0.0003) & \textbf{飽和} 3194, 0.0209 (0.0209) & \textbf{飽和}
3270, 0.0204 (0.0204) & 床 (0.0004) \\
15 & 床 (0.0003) & 床 (0.0003) & 床 (0.0003) & 床 (0.0004) \\
16 & 床 (0.0003) & \textbf{飽和} 3186, 0.0215 (0.0215) & \textbf{飽和}
3266, 0.0210 (0.0210) & 床 (0.0005) \\
\end{longtable}
}

\pandocbounded{\includegraphics[keepaspectratio,alt={図 4}]{figs/fig3_Hperp_grid_N3_N16.png}}

\textbf{図 4.} 全 54 走行の \(H_\perp/H\)（対数、40000
step、\(L=124\)、種なし、直接読出し）。黒：mp、青：hm、橙：ne、緑：rb。凡例の数値は走行前に固定した予測
\(\lambda_f\)。飽和する走行は \(10^{-27}\)
台から一直線に立ち上がり、床の走行は \(10^{-20}\)〜\(10^{-14}\)
に留まる。

\subsubsection{7.2
予測との照合（数値事実）}\label{ux4e88ux6e2cux3068ux306eux7167ux5408ux6570ux5024ux4e8bux5b9f}

\begin{itemize}
\tightlist
\item
  走行前の分類予測は \textbf{53/54}
  で実測と一致した（主結果）。唯一の不一致 rb\(_{N=13}\) は
  \(\rho-1=9.86\times10^{-4}\) で閾値 \(10^{-3}\)
  を僅かに下回り「中立」と予測したが、\(t_{50}=30809\)
  で飽和した。\(\lambda\) は予測 0.00197 に対し実測 0.00197（比
  0.998）で成長率の予測は正しい。判定規則を「予測 \(t_{50}\le\)
  走行長」に置き換えれば 54/54
  になるが、これは結果を見た後の規則変更なので主結果にはしない。
\item
  飽和した 25 走行の \(\lambda\) 実測／予測比は
  \textbf{0.997〜1.008}（指数窓の
  \(R^2\ge0.997\)）。この一致の意味は、同じ離散写像の線形化（共回転 1
  刻み線形化行列
  \(G\)）が、その写像の非線形走行の初期指数率を与えること------線形安定性解析と直接時間発展の一致------であり、独立な物理的予言ではない。
\item
  床の走行の後半（30000〜40000 step）の \(\ln f\) の傾きは
  \(2.7\)〜\(3.2\times10^{-4}\)/step で、生成法・\(N\)
  に依らない共通値である（§8.3）。mp の
  \(N=8,9\)（\(5.4\)、\(6.6\times10^{-4}\)）はこれに流れの弱い不安定性が加わり、40000
  step で \(10^{-10}\)、\(4\times10^{-9}\) まで成長した。
\end{itemize}

\pandocbounded{\includegraphics[keepaspectratio,alt={図 5}]{figs/fig5_rho_bands.png}}

\textbf{図 5.} 走行前に固定した \(\rho-1\)（共回転 1
刻み線形化行列）と実測分類。塗り＝飽和／成長中、白抜き＝床。緑の帯（\(\rho-1<10^{-3}\)）が「40000
step 内は床」の予測域。帯の境界に乗った 1 点（rb,
\(N=13\)）だけが分類を外した。

\pandocbounded{\includegraphics[keepaspectratio,alt={図 6}]{figs/fig6_lambda_pred_vs_meas.png}}

\textbf{図 6.} 成長率：予測 \(\lambda_f=2\ln\rho\) 対 実測
\(\lambda\)（\(\ln H_\perp\) の指数窓の傾き）。飽和 25 走行で比
0.997〜1.008。

\pandocbounded{\includegraphics[keepaspectratio,alt={図 7}]{figs/fig7_t50_vs_N.png}}

\textbf{図 7.} \(t_{50}\)（\(H_\perp/H=0.5\) 到達 step）の \(N\)
依存。塗り＝実測、白抜き＝予測（\(f_0=3\times10^{-32}\)）。破線は走行長
40000。hm/ne の偶数 \(N\) は \(N\) とともに \(t_{50}\) が短くなり、rb は
\(N\) とともに長くなって \(N=14\) 以上で走行長を超える。

\subsubsection{7.3
飽和後の状態（数値事実）}\label{ux98fdux548cux5f8cux306eux72b6ux614bux6570ux5024ux4e8bux5b9f}

飽和した 25 走行すべてで、step 40000 の参加率
PR/M（\(=(\sum|z|^2)^2/(M\sum|z|^4)\)）は 0.05〜0.31 で \(N\)
とともに減少し（\(N=5\)：0.27〜0.31、\(N=16\)：0.05〜0.07）、最小振幅は
\(10^{-4}\) 以下、最大振幅は
0.6〜1.05------少数の関係への局在である。等モジュラー親（mp,
hm）、非等モジュラー親（ne）、乱数均衡親（rb）のどれから出発しても同じ型に落ちる。v1
§18--19 の等分配（PR/M \(=1\)、\(|z|^2\to1/M\)）は 54
走行のどこにも現れない。

\pandocbounded{\includegraphics[keepaspectratio,alt={図 8}]{figs/fig4_final_state.png}}

\textbf{図 8.} step 40000 の状態。左：飽和した走行の参加率
PR/M------破線の PR/M \(=1\)（旧力学の等分配）に達するものはなく、\(N\)
とともに 0.3 から 0.05 へ下がる。右：振幅の広がり
std\(|z|\)/rms\(|z|\)------飽和した走行は生成法によらず 0.7〜0.95
に集まり、床の走行は初期値（等モジュラーなら 0）のまま。

\subsubsection{7.4
生成法ごとの読み（数値事実の整理。解釈は加えない）}\label{ux751fux6210ux6cd5ux3054ux3068ux306eux8aadux307fux6570ux5024ux4e8bux5b9fux306eux6574ux7406ux89e3ux91c8ux306fux52a0ux3048ux306aux3044}

\begin{enumerate}
\def\labelenumi{\arabic{enumi}.}
\tightlist
\item
  \textbf{hm と ne は同じ縞}：偶数 \(N=6,8,\dots,16\) は全て飽和、奇数
  \(N\) と \(N=3,4\) は全て床。ne は hm
  と位相配置が同じで振幅だけ異なる（\(|z|\) の幅は
  0.11〜0.42）ので、\textbf{振幅分布は安定性の帰属（飽和／床）を変えず、\(\lambda\)
  を 2〜3\% 変えるだけ}である。
\item
  偶数 \(N\) の hm/ne の \(\lambda\) は \(N\)
  とともに単調に増える（0.016 → 0.0215、\(t_{50}\) は 4338 → 3186）。
\item
  \textbf{高対称系列（hm,
  ne）には明瞭な奇偶非対称が存在する}。ただしそれは \(N\)
  の偶奇だけで決まる普遍則ではない：同じ \(N=5,7,9,11,13\)
  で対称性のない rb は飽和する。この縞は「1-因子分解（偶数 \(N\)
  にしか存在しない）／距離クラス（奇数 \(N\)
  に用いた）」という構成法の切り替わりと同時に起きているので、原因が
  \(N\) の数論的 parity
  そのものか、それに伴う離散対称構造の違いかは\textbf{未分離}である（証明済みは「偶奇だけでは説明できない」まで。偶数
  \(N\) に距離クラス構成を、奇数 \(N\) に
  1-因子分解に代わる別構成を与えて走らせれば分離できるが、本稿では未実施）。
\item
  mp（v1 系列の生成法）は \(N=5,6,7\) のみ飽和、\(N\ge8\) は床。v1
  系列で「\(N\) 依存」に見えていたものは make\_parent
  生成法で得られる親系列の性質である。
\item
  rb は \(N=5\)〜13 で飽和、\(N=14,15,16\) で床（各 \(N\) 1 個）。
\item
  \(N=3,4\) は 4 生成法とも床（\(\rho-1\le10^{-5}\)）。
\end{enumerate}

以上から、\textbf{修正後の力学でどの \(N\)
でインフレーション的発展が起きるかは、\(N\)
単独では決まらず、初期状態の構造が必要である}。自己無撞着性は状態の許容集合を切り出すが、その中の状態を一意に選択しない（非等モジュラー解、連続族、安定性の異なる親が同じ
\(N\) に共存する）。

\begin{center}\rule{0.5\linewidth}{0.5pt}\end{center}

\subsection{8.
弾道則・読出し・刻み由来の成長}\label{ux5f3eux9053ux5247ux8aadux51faux3057ux523bux307fux7531ux6765ux306eux6210ux9577}

\subsubsection{8.1 弾道則（数値事実、多倍長 50
桁）}\label{ux5f3eux9053ux5247ux6570ux5024ux4e8bux5b9fux591aux500dux9577-50-ux6841}

親の自己無撞着残差を \(r=\|iKv-\mu v\|/\|v\|\)
とすると、修正後の力学の立ち上がりは短時間で

\[
\frac{H_\perp}{H_{\rm total}}\simeq(r\tau)^2\qquad(\tau\to0),\qquad \tau=\Delta\cdot\text{step}
\]

に従う（数値事実。一次展開からの解析導出は行っていない）。mpmath 50 桁で
\(\Delta=2\pi/10^6\)・5 step
を走らせた結果（\texttt{Nall\_linear1000000\_steps5\_mpmath50\_…\_20260828}）：

{\def\LTcaptype{none} % do not increment counter
\begin{longtable}[]{@{}
  >{\raggedright\arraybackslash}p{(\linewidth - 10\tabcolsep) * \real{0.1667}}
  >{\raggedright\arraybackslash}p{(\linewidth - 10\tabcolsep) * \real{0.1667}}
  >{\raggedright\arraybackslash}p{(\linewidth - 10\tabcolsep) * \real{0.1667}}
  >{\raggedright\arraybackslash}p{(\linewidth - 10\tabcolsep) * \real{0.1667}}
  >{\raggedright\arraybackslash}p{(\linewidth - 10\tabcolsep) * \real{0.1667}}
  >{\raggedright\arraybackslash}p{(\linewidth - 10\tabcolsep) * \real{0.1667}}@{}}
\toprule\noalign{}
\begin{minipage}[b]{\linewidth}\raggedright
\(N\)
\end{minipage} & \begin{minipage}[b]{\linewidth}\raggedright
\(r\)
\end{minipage} & \begin{minipage}[b]{\linewidth}\raggedright
step 1 実測
\end{minipage} & \begin{minipage}[b]{\linewidth}\raggedright
\((r\Delta)^2\)
\end{minipage} & \begin{minipage}[b]{\linewidth}\raggedright
step 5 実測
\end{minipage} & \begin{minipage}[b]{\linewidth}\raggedright
\((r\tau)^2\)
\end{minipage} \\
\midrule\noalign{}
\endhead
\bottomrule\noalign{}
\endlastfoot
5 & \(3.78\times10^{-11}\) & \(5.63248\times10^{-32}\) &
\(5.63248\times10^{-32}\) & \(1.40812\times10^{-30}\) &
\(1.40812\times10^{-30}\) \\
8 & \(2.01\times10^{-11}\) & \(1.59347\times10^{-32}\) &
\(1.59347\times10^{-32}\) & \(3.98367\times10^{-31}\) &
\(3.98368\times10^{-31}\) \\
16 & \(9.38\times10^{-13}\) & \(3.47645\times10^{-35}\) &
\(3.47645\times10^{-35}\) & \(8.69113\times10^{-34}\) &
\(8.69114\times10^{-34}\) \\
20 & \(1.38\times10^{-12}\) & \(7.56157\times10^{-35}\) &
\(7.56245\times10^{-35}\) & \(1.89039\times10^{-33}\) &
\(1.89061\times10^{-33}\) \\
\end{longtable}
}

\(H_{\rm total}\) のドリフトは \(10^{-50}\) 台（厳密保存）。

この結果の意味は次のとおりである。\textbf{出発の底は物理が決めているのではない。}
倍精度では \(H_\perp/H\) の床が \(10^{-32}\)（振幅で
\(10^{-16}\)）に見えるが、精度を 50 桁に上げると底は \(10^{-35}\)
以下まで深くなり、そこでも同じ法則 \((r\tau)^2\) が 5〜7
桁で成り立つ。数値的な底を決めているのは (i) 計算精度と (ii)
親の自己無撞着残差 \(r\) であり、親を高精度で \(r\to0\)
まで磨けば底は際限なく深くなる（厳密な固定点からは、不安定モードが精度雑音を増幅するまで何も起きない）。したがって
v1 記事図 1 の「\(10^{-32}\) から 31 桁」も本稿 §7 の「\(10^{-24}\) から
23
桁」も、\textbf{計算精度が見せている桁数}であって物理の桁数ではなく、これらの値から物理的な初期摂動の振幅を読み取ることはできない（摂動源は未同定、§10）。桁数が
31 から 23 に減ったのではない（\(10^{-24}\) は倍精度の床 \(10^{-32}\)
に親残差 \(r\approx10^{-11}\) の弾道項 \((r\Delta)^2\)
が乗った値である）。

\pandocbounded{\includegraphics[keepaspectratio,alt={図 9}]{figs/figB_ballistic_mp50.png}}

\textbf{図 9.} 多倍長 50 桁、\(\Delta=2\pi/10^6\)、5 step の
\(H_\perp/H\)（直接読出し）。倍精度の床 \(10^{-31}\) の下でも
\(N=5,8,10,16,20\) の全てが
\((r\tau)^2\)（破線、実測と重なって見えない）に乗る。底は精度と親残差で決まり、精度を上げれば深くなる。

v1 の onset--残差則 \(t_{\rm onset}=11.616[-\ln\varepsilon]-99.6\)
は位相のみ \(K\)
の固定点についての法則であり、修正後の力学では、親残差が作る決定論的ドリフト
\(\simeq(r\tau)^2\) の上に、親の不安定性による指数成長
\(e^{\lambda_f\,\text{step}}\) が乗る形に置き換わる。\(N=5\)（mp）では
\(10^{-24}\) から \(10^{-1}\) まで 23
桁が一直線（\(\lambda=0.0089\)/step、\(R^2=1.000\)）である。

\pandocbounded{\includegraphics[keepaspectratio,alt={図 10}]{figs/figD_N5_directHperp_40000.png}}

\textbf{図 10.} \(N=5\)、make\_parent 等モジュラー親、40000
step、直接読出しの \(H_\perp\)（対数）。step 1 で倍精度の床 \(10^{-32}\)
から弾道項 \((r\Delta)^2\approx10^{-24}\) に跳び、以後 \(10^{-1}\)
まで一定率 0.0089/step の指数成長、約 5000 step
で飽和して局在状態に留まる。潜伏と見える区間の中身は最初から同じ率の指数成長である。

\subsubsection{8.2
読出し（数値事実）}\label{ux8aadux51faux3057ux6570ux5024ux4e8bux5b9f}

\(H_\perp\) を \(H_{\rm total}-H_\parallel\)
の差し引きで求めると丸め床が \(10^{-15}\)（\(H\)
の尺度）になり、それ以下の成長が「不動」に見える。\(N=16\)
を「安定」と読んだ中間記録はこの読出し不良によるもので、直交成分の直接計算（床
\(10^{-32}\)）では \(N=16\) の mp 親も 40000 step で
\(10^{-32}\to1.6\times10^{-15}\)
と単調に成長している。本稿の数値は全て直接読出しである。

\subsubsection{8.3
刻み由来の成長（数値事実）}\label{ux523bux307fux7531ux6765ux306eux6210ux9577ux6570ux5024ux4e8bux5b9f}

刻み角 \(\Delta\) を \(L=62\)〜1984
で変えた走行（\texttt{飽和ステップ数とNの関係\_固定点ヤコビアン解析\_20260829}）から、\(H_\perp/H\)
の 1 刻み成長率は

\[
\lambda_f=a\,\Delta+b\,\Delta^2,\qquad b=0.1250\simeq\tfrac18\ (N=16,\ 4\text{ 桁の数値結果})
\]

に分離される。\(a\) は連続流 \(dv/d\tau=K(v)v\)
の相対平衡の線形不安定性（離散化に依存しない模型内の不安定性。親依存で、mp
親では \(N=6,7\) で全親 \(0.13\)〜\(0.30\)、\(N\ge11\) で 25 親中
0）、\(b\Delta^2\) は積分器が状態依存の \(K\) を 1
刻み凍結することによる離散化誤差で \(N\)・親に依らない。§7
の床の後半傾き \(\sim3\times10^{-4}\)/step はこの第 2
項（\(\Delta=2\pi/124\) で \(\Delta^2/8=3.2\times10^{-4}\)）である。

\pandocbounded{\includegraphics[keepaspectratio,alt={図 11}]{figs/figC_lambda_t50_vs_N.png}}

\textbf{図 11.} make\_parent 等モジュラー親（各 \(N\) 5
実現）の成長率と飽和 step。左：共回転 1 刻み線形化行列の
\(2\lambda_G\)（青、中央値と実現ごとの点）と実測
\(\lambda\)（赤四角）------\(N\ge9\) で \(3\times10^{-4}\)/step
の共通床（刻み由来項）に張り付く。右：予測
\(t_{50}=(\ln0.5-\ln f_{\rm seed})/\lambda_f\) と実測------\(N=5\)〜7 は
40000 以下、\(N\ge8\) は \(10^5\) 以上。\(\tau\) 単位では
\(\tau_{50}=355/\Delta\to\infty\)（\(\Delta\to0\)）なので、\textbf{中立な親は連続極限では飽和しない}。§7
の「床」は 40000 step 内の判定であり、刻み由来の成長は \(10^5\) step
台で \(O(1)\) に達しうる。

\begin{center}\rule{0.5\linewidth}{0.5pt}\end{center}

\subsection{9. v1
の主張の維持／修正／撤回}\label{v1-ux306eux4e3bux5f35ux306eux7dadux6301ux4feeux6b63ux64a4ux56de}

{\def\LTcaptype{none} % do not increment counter
\begin{longtable}[]{@{}
  >{\raggedright\arraybackslash}p{(\linewidth - 4\tabcolsep) * \real{0.3333}}
  >{\raggedright\arraybackslash}p{(\linewidth - 4\tabcolsep) * \real{0.3333}}
  >{\raggedright\arraybackslash}p{(\linewidth - 4\tabcolsep) * \real{0.3333}}@{}}
\toprule\noalign{}
\begin{minipage}[b]{\linewidth}\raggedright
v1 の主張（§28 の番号、記事図）
\end{minipage} & \begin{minipage}[b]{\linewidth}\raggedright
判定
\end{minipage} & \begin{minipage}[b]{\linewidth}\raggedright
根拠
\end{minipage} \\
\midrule\noalign{}
\endhead
\bottomrule\noalign{}
\endlastfoot
1. 実直交更新により \(Z^\dagger Z\)・\(Z^TZ\) は厳密保存 &
\textbf{維持}（Cayley→\(\exp\) に書換） & §2.2、§8.1（\(10^{-50}\)） \\
2. \(H_\perp\) の増大は内部移送、\(H_\perp\le H_{\rm total}\) &
\textbf{維持} & §7 \\
3. \(K/\sigma_{\max}\) 正規化は時計の付け替え、6.8\%
差は刻み収束則で説明 & \textbf{限定} & 位相のみ \(K\)
についての記述。修正後は \(K/\sigma\) 枝を廃止 \\
4. onset \(\propto-\ln\varepsilon\)、成長率不変 & \textbf{修正} &
§8.1：弾道則 \((r\tau)^2\)＋親依存の \(\lambda_f\) \\
5. Floquet 実 2 重固有値 \(\mu_1=1.0901\)、rank-4 選択 & \textbf{修正} &
旧写像の値。修正後の \(N=5\) mp 親は共回転 1 刻み線形化行列で
\(\rho=1.0045\)、hm 親は中立 \\
6. 三重整合 0.1725／11.616／\(\mu_1\) & \textbf{修正} & §7.2：1
刻み線形化行列の予測と実測 \(\lambda\) の一致（0.997〜1.008）に置換 \\
7. star 閉塞 \(\Leftrightarrow\) 頂点が複素ヌル錐上 &
\textbf{維持}（用語を複素ヌル錐に） & §5.3 \\
8. 準安定状態＝等モジュラー・ヌル複素シンプレックス（等分配） &
\textbf{撤回} & §7.3：修正後は局在。等分配は位相のみ \(K\) の終着点 \\
9. \(N=4\) の \(120°\) は局所閉塞＋等分配から定理 &
\textbf{維持（定理）／撤回（到達）} &
定理は正しい。系がそこへ到達するのは旧力学 \\
10. \(N=5\) の 13 個・12 通り閉包 & \textbf{撤回} &
旧力学でも軌道依存（再実行で 0）、修正後は 4 群構造が現れない \\
11. 8 seed で \(3+3+2+2\) と等モジュラス & \textbf{撤回} & 同上 \\
12. 自己無撞着 ⇒ 閉塞 ⇒ \(S^1\) 軌道、\(U^n=I\) は別途 & \textbf{維持} &
§4 \\
記事図 1「\(10^{-32}\) から 31 桁」 & \textbf{修正（強める方向）} &
出発の底 \(10^{-32}\) は倍精度の丸め床であって物理の底ではない。多倍長
50 桁で同じ走行を行うと底はさらに深くなり、\(10^{-35}\) まで
\((r\tau)^2\) の法則が続く（§8.1）。数値的な底は計算精度と親残差 \(r\)
で決まる下限であり、\(r\to0\) の厳密な固定点では任意に深い。「31
桁」「23 桁」は計算精度が見せる桁数であって、物理の桁数が 31 から 23
に減ったのではなく、これらから物理的な初期摂動振幅は読み取れない \\
記事図 6「完全均等分配」 & \textbf{撤回} & §7.3 \\
§20 の \(N=3\)〜16 表（rank \(=N-1\)、\(|z|^2\to1/M\)） & \textbf{修正}
& rank は維持。等分配は撤回。\(\sigma=N-1\) は \(\mu=-(N-1)r^2\)
として維持（\(N=3\) は \(-\tfrac32r^2\)） \\
§4「\texttt{make\_parent} は監査済み」 & \textbf{訂正} & §1.1、§3.1 \\
§25.3--25.4 の IIB 行列模型・動的コンパクト化との類似、§26.3 の
preheating 型類似 & \textbf{削除} & 図 1・6
に依拠していた考察。本稿は新解釈を入れない \\
\end{longtable}
}

\begin{center}\rule{0.5\linewidth}{0.5pt}\end{center}

\subsection{10. 限界}\label{ux9650ux754c}

\begin{enumerate}
\def\labelenumi{\arabic{enumi}.}
\tightlist
\item
  \textbf{軌道の計算機依存}：長時間の個別軌道は丸め誤差の指数増幅で機械間で再現しない（\(N=5\)
  で 6800 step 以降、\(10^{-8}\to10^{-2}\) に 2600 step）。本稿の
  \(t_{50}\)・\(\lambda\)・PR/M
  は統計量として安定だが、飽和後の個別配置は再現しない。v1 の \(N=5\)
  主張（\(3+3+2+2\)、13/12）は同じ旧力学でも再実行で現れなかった。
\item
  \textbf{rb は各 \(N\) 1 個}：\(N\ge14\) の中立は 1
  例ずつであり、統計ではない。
\item
  \textbf{床は 40000 step 内の判定}：§8.3 の刻み由来項により、中立な親も
  step 単位では \(10^5\) 台で \(O(1)\) に達しうる。連続極限の主張には
  \(\tau\) 単位の不変性確認が要る。
\item
  \textbf{\(H_{\rm int}\) を厳密保存しない}：\(\exp(\Delta K(z))\) は
  \(K\) を 1 刻み凍結するため
  \(H_{\rm int}=\tfrac12\sum_{\{e,f\}}(\mathrm{Im}\,\bar z_ez_f)^2\) が
  \(\Delta\) に比例して漂う（\(\Delta\) を 1/10 にすると変動も
  1/10）。ノルムと閉塞は厳密保存される。積分器がシンプレクティックか否かは本稿では判定していない（判定には
  1 刻み写像の実ヤコビアンについて \((D\Phi)^TJ\,D\Phi-J\)
  を調べる必要がある）。
\item
  \textbf{未証明命題}：(a) \(N\ge5\) で「自己無撞着 \(\Rightarrow\)
  局所閉塞」（等モジュラー点近傍では数値的に成立、\(N=4\)
  には反例分枝）。(b) 対称な親（hm の奇数
  \(N\)、\(N=3,4\)）が中立になる機構。(c) 高対称系列の奇偶非対称の原因が
  parity か構成法の離散対称構造か（§7.4）。(d) \(U^n=I\)（有理ロック）。
\item
  \textbf{摂動源は未同定}：厳密な自己無撞着状態は
  \(v(\tau)=e^{-i\mu\tau}v(0)\) のまま離れない。本稿の立ち上がりは親残差
  \(r\)
  と丸め誤差が不安定モードに入って始まっており（§8.1）、本稿が確認したのは線形不安定性とその非線形飽和であって、物理的な摂動源を含む開始機構ではない。
\item
  \textbf{fixed 系の \(N=6,7,10,11\)} は親探索の汚染で対照無効（§3.3）。
\end{enumerate}

\begin{center}\rule{0.5\linewidth}{0.5pt}\end{center}

\subsection{11. 結論}\label{ux7d50ux8ad6}

\begin{enumerate}
\def\labelenumi{\arabic{enumi}.}
\tightlist
\item
  ゼロ閉塞 \(\sum z^2=0\)
  は、生成子が実反対称で状態がその非零固有モードであることだけから従う定理であり、正規化・振幅分布・\(N\)・親の生成法に依らない（証明済み。110
  親の実装で数値検証）。
\item
  複素シンプレックスは状態空間を 1 点も削らない。任意の複素 2 乗距離は
  Takagi 分解で厳密に埋め込まれ（証明済み。乱数 1400
  状態で数値検証）、閉塞が可能なのは複素化が正値性を除去したからである。形が表すのは状態の符号を忘れた像
  \(v/(\mathbb Z_2)^M\)
  で、符号枝は力学の厳密な離散対称性で結ばれる。幾何は結果を描く言語であって状態を選ぶ原理ではない。
\item
  隠れた振幅正規化の除去・凍結生成子の指数回転・厳密な自己無撞着初期化のもとでも、自己無撞着な相対平衡のうち線形不安定な親に微小摂動（親残差・丸め誤差）が入ると横方向成分は指数成長し非線形飽和する。成長率は
  v1 の 10〜30 分の 1、飽和後は等分配でなく局在である。高対称な hm/ne
  系列では \(N=6\)〜16 の偶数と \(N=5\)〜15
  の奇数の間に明瞭な安定性の非対称が現れたが（\(N=3,4\)
  は床）、乱数均衡親では奇数 \(N\) にも不安定状態が存在するので、これは
  \(N\) の偶奇だけで決まる普遍則ではない。したがって \(N\)
  単独は発展の有無を決めず、初期状態の構造が必要である。離散写像の共回転
  1 刻み線形化行列は 54 走行の帰属を 53/54、成長率を 0.997〜1.008
  で与える。
\end{enumerate}

まとめると：閉塞は定理として残った。複素シンプレックスは選択原理ではなかった。旧インフレーション像（潜伏→急拡大→等分配）は修正されたが、連続流の相対平衡の線形不安定性とその飽和は残った。高対称系列には奇偶非対称が現れた。自己無撞着性だけでは物理状態を一意に選択できず、追加の選択原理は未発見である。

\begin{center}\rule{0.5\linewidth}{0.5pt}\end{center}

\subsection{12.
引用と再現性}\label{ux5f15ux7528ux3068ux518dux73feux6027}

\subsubsection{12.1 自己引用}\label{ux81eaux5df1ux5f15ux7528}

\textbf{{[}K1{]}} N. Kihara, \emph{Anonymous Equal-Amplitude
Composite-Wave Model: Basic Axiom System v9 --- Pure Definition}, 2026.
\textbf{{[}K2{]}} 木原範昭,
「N体完全二体関係波における生成子ランク線形上界と三方向飽和」, 2026.
Concept DOI: 10.5281/zenodo.21465898. \textbf{{[}K8{]}} 木原範昭,
「N体関係波閉鎖系における三方向生成の時間構造――二段階seed除去による因果分離」,
2026. Concept DOI: 10.5281/zenodo.21614402. \textbf{{[}K9{]}} 木原範昭,
「幾何級数的急拡大は不安定な自己無撞着閉包に固有である------一般零閉塞初期状態による開始様式の因果判別」,
2026. Concept DOI: 10.5281/zenodo.21798854. \textbf{{[}v1{]}} 木原範昭,
本論文第 1 版, 2026-08-27. Version DOI: 10.5281/zenodo.22112009.

\subsubsection{12.2 外部引用}\label{ux5916ux90e8ux5f15ux7528}

\textbf{{[}E1{]}} A. Aste, ``Origin of the Complex Structure of Quantum
Mechanics,'' arXiv:1905.12894, 2019. \textbf{{[}E2{]}} T. Takagi, ``On
an algebraic problem related to an analytic theorem of Carathéodory and
Fejér and on an allied theorem of Landau,'' \emph{Japan. J. Math.} 1
(1925) 83--93. （Autonne--Takagi 分解。R. A. Horn, C. R. Johnson,
\emph{Matrix Analysis}, 2nd ed., Cambridge Univ. Press, 2013, Cor. 4.4.4
も参照） \textbf{{[}E3{]}} I. J. Schoenberg, ``Remarks to Maurice
Fréchet's article `Sur la définition axiomatique d'une classe d'espace
distanciés vectoriellement applicable sur l'espace de Hilbert',''
\emph{Ann. of Math.} 36 (1935) 724--732.

\subsubsection{12.3
再現パッケージ}\label{ux518dux73feux30d1ux30c3ux30b1ux30fcux30b8}

v1 が引用した 15 パッケージ（旧エンジン。うち 14 zip は Zenodo 22112009
に収録、記事図 1 の源
\texttt{complex\_simplex\_decompactification\_N5\_N16\_20260826}
は本リポジトリ収録）は v1
の再現用としてそのまま残す。本稿の主張は以下のパッケージ（GitHub
\texttt{WurabeSeiji/ai-chat-logs-open}、\texttt{次元の生成構造/電子の反跳実験/seed除去による準安定相追試/chatgpt追試/}
配下、各パッケージに
program・data・results・figures・README・run\_all.sh・SHA256SUMS
を同梱）で再現できる。

{\def\LTcaptype{none} % do not increment counter
\begin{longtable}[]{@{}
  >{\raggedright\arraybackslash}p{(\linewidth - 2\tabcolsep) * \real{0.5000}}
  >{\raggedright\arraybackslash}p{(\linewidth - 2\tabcolsep) * \real{0.5000}}@{}}
\toprule\noalign{}
\begin{minipage}[b]{\linewidth}\raggedright
節
\end{minipage} & \begin{minipage}[b]{\linewidth}\raggedright
パッケージ
\end{minipage} \\
\midrule\noalign{}
\endhead
\bottomrule\noalign{}
\endlastfoot
§3.2 & \texttt{論文v1\_全再現テスト\_20260828} \\
§3.3 &
\texttt{論文v1\_全プログラム修正版\_20260828}（fixed／fixed\_baseline／fixed\_equimodular、四者比較）、\texttt{不足プログラム復元\_chatgpt\_20260828} \\
§4 &
\texttt{公理見直し\_ゼロ閉塞定理と固有時計\_20260829}、\texttt{v2補完実験\_4生成法\_N3toN16\_統一プロトコル\_20260830} \\
§5 &
\texttt{v2補完実験\_…}（pass2）、\texttt{複素シンプレックス基礎\_N別全展開\_20260830}、\texttt{複素シンプレックス基礎\_N別全展開\_非等モジュラー版\_20260830} \\
§6 &
\texttt{手作り自己無撞着親と対称性\_倍音と関係数の検討\_20260829}、\texttt{複素シンプレックス\_重心閉塞\_非等モジュラー族\_20260830} \\
§7 & \texttt{v2補完実験\_4生成法\_N3toN16\_統一プロトコル\_20260830}（54
走行、予測、図 2〜8） \\
§8 &
\texttt{N5/N16\_linear124\_equimodular\_selfconsistent\_directHperp\_treatment\_only\_20260828}、\texttt{Nall\_linear1000000\_steps5\_mpmath50\_…\_20260828}、\texttt{飽和ステップ数とNの関係\_固定点ヤコビアン解析\_20260829} \\
\end{longtable}
}

\subsubsection{12.4
図ファイルの配置}\label{ux56f3ux30d5ux30a1ux30a4ux30ebux306eux914dux7f6e}

本文の図は \texttt{figs/}（本文と同じフォルダ）に集約した。出所：

\begin{itemize}
\tightlist
\item
  図 1 \texttt{figA\_four\_way\_Hperp\_N3\_N16.png} ←
  \texttt{論文v1\_全プログラム修正版\_20260828/results/figures/four\_way\_Hperp\_frac\_N3\_N16.png}
\item
  図 2 \texttt{fig1\_closure\_step0.png}、図 3
  \texttt{fig2\_embed\_random.png}、図 4
  \texttt{fig3\_Hperp\_grid\_N3\_N16.png}、図 5
  \texttt{fig5\_rho\_bands.png}、図 6
  \texttt{fig6\_lambda\_pred\_vs\_meas.png}、図 7
  \texttt{fig7\_t50\_vs\_N.png}、図 8 \texttt{fig4\_final\_state.png} ←
  \texttt{v2補完実験\_4生成法\_N3toN16\_統一プロトコル\_20260830/figures/}
\item
  図 9 \texttt{figB\_ballistic\_mp50.png} ←
  \texttt{Nall\_linear1000000\_steps5\_mpmath50\_…\_20260828/figures/compare\_N\_L1000000\_5\_mp50.png}
\item
  図 10 \texttt{figD\_N5\_directHperp\_40000.png} ←
  \texttt{N5\_linear124\_equimodular\_selfconsistent\_directHperp\_treatment\_only\_20260828/figures/N5\_Hperp\_baseline\_vs\_amplitude\_aware.png}（ファイル名は旧スクリプト固定、中身は
  treatment のみ）
\item
  図 11 \texttt{figC\_lambda\_t50\_vs\_N.png} ←
  \texttt{飽和ステップ数とNの関係\_固定点ヤコビアン解析\_20260829/figures/lambda\_and\_t50\_vs\_N.png}
\end{itemize}

\end{document}
